\documentclass[11pt]{article}

\usepackage[margin=1.1in]{geometry}
\usepackage{amsmath,amssymb}
\usepackage{booktabs}
\usepackage[round]{natbib}
\usepackage{microtype}
\usepackage{url}
\usepackage[hidelinks]{hyperref}
\usepackage{setspace}

\newcommand{\brier}{\mathrm{Brier}}

\title{Transparent Base Rates Approach the State of the Art\\in Conflict Forecasting\\[6pt]
\large Evidence from a fully auditable forecasting system, five preregistered\\covariate failures, and equal-weight forecast pooling}

\author{Kara Zajac\thanks{Independent researcher. Email:
\texttt{kara@soulstone.org}. Code (MIT) and the merged dataset (CC~BY~4.0)
are public at \url{https://github.com/KaraZajac/TOCSIN}; a continuously
updating deployment, including all diagnostics reported here, is at
\url{https://tocsin.karazajac.io}. A single command regenerates every number
in this paper from primary sources.}}

\date{Preprint \,\textperiodcentered\, July 2026 \,\textperiodcentered\, v1.0}

\begin{document}
\maketitle

\begin{abstract}
\noindent Machine-learning systems now define the state of the art in
armed-conflict forecasting, but their opacity makes them difficult to audit,
and a recurring finding---that naive baselines are nearly as accurate---has
never been given a systematic, adversarially honest treatment by a system
\emph{designed} to be transparent. We present TOCSIN, a conflict-forecasting
system whose every probability is a reference-class base rate: an empirically
shrunk historical frequency, conditioned on recency structure, auditable down
to the individual class counts. Evaluated walk-forward on the same target,
vantages, and information sets as the Uppsala VIEWS system (7{,}680
country-months across five historical vantages), TOCSIN attains a Brier score
of 0.0461 against VIEWS's 0.0412---within 12\% of the state of the art, while
beating it on 70\% of individual country-months---and naive last-year
persistence (0.0433) falls between them. We then subject the system to a
preregistered tune/validate protocol and test five families of theoretically
motivated covariates: regime type, forecast-horizon decay, youth bulges,
militarized-dispute and alliance history, and ethnic exclusion (alone and
interacted). All five fail out-of-sample, despite large descriptive
associations (youth bulges quadruple onset rates; militarized-dispute history
marks a 30-fold onset class among long-quiet country pairs). We identify
three mechanisms---absorption by recency conditioning, class fragmentation
under shrinkage, and, at rare-event base rates, a formal \emph{metric
blindness} of the Brier score, for which we give an exact decomposition
showing that even perfect conditioning cannot move aggregate Brier by more
than 0.05\% at pair-grain base rates. Finally, zero-parameter equal-weight
pooling of VIEWS with persistence outperforms every individual model (0.0399,
$-3.3\%$ versus VIEWS), while adding TOCSIN to that pool \emph{worsens} it:
the transparent system wins months but duplicates the pool's recency
information rather than diversifying it. The results sharpen a deflationary
view of conflict predictability---the practical ceiling on the standard
occurrence target sits close to what recency alone delivers, and is most
cheaply approached not by richer models but by pooling diverse ones.\\[4pt]
\noindent\textbf{Keywords:} conflict forecasting; base rates;
reference-class forecasting; Brier score; forecast combination;
preregistration; UCDP; VIEWS
\end{abstract}

\section{Introduction}\label{sec:intro}

Armed-conflict research has taken a decisive predictive turn. Where
\citet{ward2010perils} diagnosed a field optimizing statistical significance
at the expense of predictive power, there now exist standing forecast
systems---most prominently the Uppsala Violence Early-Warning System
\citep[VIEWS;][]{hegre2019views}---that publish monthly machine-learning
forecasts of state-based violence for every country, and organized
competitions that score them \citep{hegre2022lessons,vesco2022united}. The
predictive turn, however, has produced a persistent embarrassment alongside
its successes: simple baselines---last year's violence, the country's
historical rate---come uncomfortably close to the most sophisticated
ensembles \citep{cederman2017predicting,chadefaux2017limits}. This gap
between methodological ambition and marginal accuracy gain is usually
reported in passing, as a caveat. It has rarely been made the \emph{object}
of study.

This paper does so, from an unusual direction. Rather than build a richer
model and report its increment, we build the most disciplined possible
version of the \emph{simple} alternative---a pure reference-class base-rate
engine in the tradition of clinical-versus-statistical prediction
\citep{meehl1954clinical,dawes1979robust} and outside-view forecasting
\citep{kahneman1993timid}---and ask three questions with it:

\begin{enumerate}
\item \textbf{How close does transparent get?} When every probability is a
shrunk historical frequency whose class definition, member counts, and
shrinkage weights are inspectable, how much accuracy is sacrificed relative
to the covariate-rich state of the art, scored on the state of the art's own
target?
\item \textbf{What do covariates actually add, out-of-sample?} The conflict
literature proposes many structural correlates of onset---regime type
\citep{goldstone2010global}, development
\citep{fearon2003ethnicity,collier2004greed}, youth bulges
\citep{urdal2006clash}, ethnic exclusion \citep{cederman2010why}, and dyadic
dispute history \citep{maoz2019dyadic}. We test five covariate families under
a preregistered tune/validate protocol designed to make post-hoc
rationalization impossible, and report all results, including---especially---
the failures.
\item \textbf{Where do the remaining gains live?} If single-model
improvements are exhausted, the forecast-combination literature
\citep{bates1969combination,clemen1989combining,timmermann2006forecast}
predicts that \emph{pooling} diverse forecasters should beat the best of
them. We test zero-parameter equal-weight pools against every individual
system, and use the pool's composition to diagnose what information each
system actually contributes.
\end{enumerate}

The answers, in brief: transparent gets within 12\% of the state of the art
(and naive persistence within 5\%); all five covariate families fail
out-of-sample despite large descriptive associations, for reasons we can
mechanically identify; and an equal-weight pool of the ML system with naive
persistence beats everything, while our own system---despite winning 70\% of
individual months---\emph{subtracts} from that pool because its information
is already inside it. Together these results locate the practical ceiling of
the standard conflict-occurrence target close to what recency structure alone
delivers, and suggest the field's cheapest remaining accuracy lies in
combination and in target design rather than in additional structure.

A secondary contribution is infrastructural. The entire system---a merged,
crosswalked, CC~BY-licensed dataset spanning nine sources; the engine; the
evaluation harness; a self-resolving forecast journal; and a public,
continuously updating deployment---is open and reproducible. The
preregistration machinery (\S\ref{sec:protocol}) is, to our knowledge, the
first covariate-adoption protocol for a conflict-forecasting system that
mechanically enforces a single consultation of held-out data and the
retirement of spent validation eras.

\section{Related work}\label{sec:related}

\paragraph{Conflict forecasting systems.} The Political Instability Task
Force established that parsimonious models with regime-type interactions
could anticipate instability onsets years ahead \citep{goldstone2010global}.
VIEWS \citep{hegre2019views,hegre2022lessons} is the current reference
system: a monthly-updated ensemble over UCDP data producing, among other
targets, the probability of ${\geq}25$ state-based battle deaths per
country-month---the target we adopt for head-to-head comparison.
\citet{blair2020forecasting} argue theory-structured models travel better
out-of-sample; \citet{muchlinski2016comparing} document large in-sample gains
from random forests that \citet{colaresi2017robot} show require careful
out-of-sample discipline to interpret.

\paragraph{The baseline embarrassment.} \citet{cederman2017predicting}
caution that conflict's rare-event structure and the dominance of
history-dependence put low ceilings on achievable gains;
\citet{chadefaux2017limits} surveys the same pattern. Our contribution is to
\emph{instrument} this claim: we reproduce it independently (persistence
within 5\% of VIEWS on 7{,}680 scored country-months), quantify the covariate
increment under preregistration ($\approx 0$), and give an exact scoring-rule
decomposition for why large class-rate signals can be invisible to aggregate
Brier at rare-event base rates.

\paragraph{Statistical versus clinical prediction and the outside view.}
The engine is a deliberate descendant of the actuarial tradition---
\citet{meehl1954clinical}, \citet{dawes1979robust}, and the reference-class
forecasting of \citet{kahneman1993timid}. Its wager is that in a domain
governed by sticky, recurrent processes, a disciplined outside view captures
most extractable signal. The five covariate failures in
\S\ref{sec:covariates} are, in effect, a measured defense of that wager.

\paragraph{Forecast combination.} That combinations of diverse forecasts
outperform their best member is among the most replicated findings in the
forecasting literature, from \citet{bates1969combination} through
\citet{clemen1989combining} and the M-competitions \citep{makridakis2000m3}.
Equal weights are notoriously hard to beat \citep[the ``forecast combination
puzzle'';][]{timmermann2006forecast}. Section~\ref{sec:pools} imports this
machinery into the conflict domain with a zero-parameter design that
forecloses fitting-based objections, and adds a compositional diagnostic:
which member's \emph{removal} improves the pool.

\paragraph{Scoring rules.} We score with the Brier score
\citep{brier1950verification}, a strictly proper rule
\citep{gneiting2007strictly}, and report skill relative to climatology.
Section~\ref{sec:blindness} formalizes a limitation of unweighted proper
scores under extreme class imbalance that connects to the literature on
rare-event verification and weighted scoring \citep{ferro2011extremal}:
propriety guarantees honest elicitation, not sensitivity to refinements of
small-mass classes.

\section{Data}\label{sec:data}

All inputs are public. The system merges nine sources onto a common
spine---Gleditsch--Ward (G-W) country code $\times$ year
\citep{gleditsch1999interstate}---with year-aware crosswalks and a validation
gate; the merged product is redistributed under CC~BY~4.0 as
\texttt{tocsin-dataset-26.1} with a full codebook.

\begin{itemize}
\item \textbf{UCDP/PRIO} (release 26.1): the Armed Conflict Dataset and its
dyadic version \citep{gleditsch2002armed,davies2026organized}, the
Georeferenced Event Dataset \citep{sundberg2013introducing} with monthly
\emph{candidate} updates, non-state and one-sided tables, and the
termination/episode coding of \citet{kreutz2010how}. Events with UCDP
placeholder actors are excluded from dyad grain; candidate files are
deduplicated by event id with latest-version precedence.
\item \textbf{Gleditsch--Ward} system membership and minimum distances
\citep{gleditsch2001measuring}, defining state-system exposure and the
${\leq}400$\,km proximity relation.
\item \textbf{PRIO Battle Deaths v3.1} \citep{lacina2005monitoring},
extending state-based death series to 1946 for descriptive use.
\item \textbf{Powell--Thyne coups} \citep{powell2011global}: 1950--present
attempts and outcomes.
\item \textbf{V-Dem Regimes of the World} \citep{luhrmann2018regimes}, via
Our World in Data.
\item \textbf{OWID population} (UN WPP/HYDE composite), carried forward to
the data edge.
\item \textbf{World Bank WDI} covariates (income, inflation, age structure,
urbanization, infant mortality; CC~BY~4.0).
\item \textbf{Ethnic Power Relations 2021} \citep{vogt2015integrating}:
politically excluded population shares.
\item \textbf{Correlates of War}: dyadic militarized interstate disputes 4.03
\citep{maoz2019dyadic,palmer2022mid5}, National Material Capabilities v7
\citep{singer1972capability}, and Formal Alliances v4.1
\citep{gibler2009international}, with a unit-tested COW${\to}$G-W crosswalk
(unified Germany $255{\to}260$, Yemen $679{\to}678$, Serbia via the
year-aware $345{\to}340$ alias, and the Pacific microstates remapped
simultaneously).
\end{itemize}

Two engineered universes matter below. The \textbf{country universe} counts
main-system G-W members (12{,}703 country-years, 1946--2025). The
\textbf{pair universe} operationalizes politically relevant pairs as
proximity (${\leq}400$\,km) $\cup$ same-region $\cup$ P5-membership: 243{,}609
undirected pair-years, within which 98.2\% of all UCDP state-based dyadic
activity occurs---constructed \emph{ex ante}, never by conditioning on
observed conflict.

\section{The engine}\label{sec:engine}

\subsection{Estimand and reference classes}

For a unit $u$ (country, dyad, or pair), year $t$, and measure $m$ (deaths
above a threshold; UCDP activity; episode termination; coup attempt), the
engine estimates $\Pr(m \text{ holds for } u \text{ in } t)$ as a shrunk
historical frequency over a reference class defined entirely by
\emph{recency structure} observable before $t$:

\begin{itemize}
\item \textbf{Active units} are banded by episode age (consecutive
hit-years: 1, 2--3, 4--9, 10+) $\times$ last-year intensity (below/above the
UCDP war line of 1{,}000 deaths; for coups, failed/successful).
\item \textbf{Non-active units} are banded by time since last hit
(\emph{recent} 2--3, \emph{dormant} 4--10, \emph{cold} 11+/never), with a
spatial-contagion flag ($+$nbr) when any ${\leq}400$\,km neighbor was in
state-based conflict the prior year.
\item At month grain, active units additionally carry a \textbf{tempo band}
(trailing hit-months 1--3, 4--8, 9--12), and windows that are not
calendar-aligned are priced on a rolling monthly substrate whose bucket
construction provably reduces to the annual engine's for January-aligned
windows.
\end{itemize}

\subsection{Estimation}

Within bucket $b$ at level $\ell \in \{\text{self}, \text{region},
\text{global}\}$, with $k$ hits in $n$ class-years, the class rate is the
Jeffreys-smoothed frequency $\hat p = (k + \tfrac12)/(n + 1)$. Levels are
combined by empirical-Bayes shrinkage \citep{efron1975data}: the self-level
posterior is
\begin{equation}
\tilde p_{\text{self}} \;=\;
\frac{k_{\text{self}} + M\,\hat p_{\text{parent}}}{n_{\text{self}} + M},
\end{equation}
with the prior strength $M$ moment-matched from the dispersion of class
rates across units and clamped to $[5, 1000]$; regions falling below 30
class-years fall back to the global class. The headline probability is the
deepest level with adequate support, and every output ships the full
ladder---counts, rates, $M$, and posterior at each level---plus caveat notes
(left-censoring at the substrate edge, window approximations,
provisional-data flags).

\subsection{Walk-forward evaluation}

All evaluation is walk-forward: scoring year $t$ uses classes built from
data strictly before $t$ (five-year burn-in), so every reported score is
out-of-sample by construction. A class-signature memoization makes the full
seven-suite backtest ($\approx$290{,}000 unit-year records) run in seconds,
which is what renders the preregistration protocol of \S\ref{sec:protocol}
practical. Table~\ref{tab:backtest} reports current walk-forward performance
(skill $= 1 - \brier/\brier_{\text{climatology}}$).

\begin{table}[t]
\centering
\caption{Walk-forward performance across the seven evaluation suites.}
\label{tab:backtest}
\begin{tabular}{lrrrr}
\toprule
suite & $n$ & base rate & Brier & skill \\
\midrule
country, sb deaths ${\geq}25$        & 5{,}547   & .167  & .0411 & $+70.6\%$ \\
country, all-type deaths ${\geq}100$ & 5{,}547   & .160  & .0366 & $+72.8\%$ \\
country, UCDP activity               & 11{,}274  & .170  & .0495 & $+65.0\%$ \\
dyad, UCDP activity                  & 24{,}876  & .115  & .0405 & $+60.6\%$ \\
dyad, episode termination            & 2{,}797   & .234  & .1668 & $+7.1\%$  \\
pair, onset/activity                 & 234{,}136 & .0006 & .0005 & $+19.6\%$ \\
country, coup attempt                & 6{,}145   & .019  & .0197 & $+17.6\%$ \\
\bottomrule
\end{tabular}
\end{table}

Two design decisions were themselves settled empirically and are kept as
negative results: horizon-aware class decay (pricing a forecast $g$ years
past the data edge by class rates offset $g$) measured \emph{worse}
($+.003$ Brier; class-level decay pools exiting with persisting units), so
one-step rates are applied frozen; and a promote-only nowcast lets
provisional current-year data upgrade but never downgrade a bucket.

\section{The arena: scoring against the state of the art}\label{sec:arena}

\subsection{Design}

We score all systems on VIEWS's native target---$\Pr({\geq}25$ state-based
battle deaths in a country-month$)$, UCDP best estimate---at each of five
historical vantages (VIEWS \texttt{fatalities002} runs of April and October
2023, April and October 2024, April 2025), horizons 1--12 months, using only
information available at each vantage. TOCSIN runs walk-forward-clamped
(classes end at the vantage). The VIEWS probability is its published
\texttt{main\_dich}, whose semantics we verified empirically rather than
assumed: across pre-2024 runs its mean (.138) matches the realized ${\geq}25$
frequency (.136), not the ${\geq}1$ frequency (.251). Baselines:
\textbf{persistence} (country's trailing-12-month hit rate, add-one smoothed)
and \textbf{climatology} (full-history monthly rate, Jeffreys-smoothed).
7{,}680 scored country-months; 7\% of outcomes are provisional (candidate
GED) at scoring time.

\subsection{Results}

\begin{table}[t]
\centering
\caption{The arena: Brier scores on $\Pr({\geq}25$ sb deaths per
country-month$)$, five vantages, horizons 1--12. Pools are equal-weight
linear opinion pools (zero fitted parameters).}
\label{tab:arena}
\begin{tabular}{lrrrrr}
\toprule
model & Brier & skill & h1--3 & h4--6 & h7--12 \\
\midrule
\textbf{pool: VIEWS + persistence} & \textbf{.03987} & $+56.8\%$ & .0340 & .0376 & .0439 \\
pool: all three                    & .04087 & $+55.7\%$ & .0354 & .0382 & .0450 \\
VIEWS                              & .04123 & $+55.3\%$ & .0337 & .0396 & .0458 \\
persistence                        & .04326 & $+53.1\%$ & .0389 & .0401 & .0471 \\
TOCSIN                             & .04610 & $+50.0\%$ & .0413 & .0425 & .0503 \\
climatology                        & .09219 & 0         & .0877 & .0929 & .0940 \\
\bottomrule
\end{tabular}
\end{table}

Table~\ref{tab:arena} reports the standings. Three readings, before the
pools (deferred to \S\ref{sec:pools}):

\begin{enumerate}
\item \textbf{The transparency tax is ${\sim}12\%$.} A pure lookup
table---auditable to the class counts---concedes .0049 Brier to a
covariate-rich ML ensemble on the ensemble's home target. Half of an
initially larger gap was closed by a single recency-structure refinement
(the tempo band moved TOCSIN from .0611 to .0461), i.e., by \emph{better
history}, not covariates.
\item \textbf{The baseline embarrassment replicates.} Persistence sits
within 5\% of VIEWS. This is an independent reproduction, on a fixed public
target with verified semantics, of the pattern \citet{cederman2017predicting}
warned about.
\item \textbf{Aggregate and disaggregate disagree.} TOCSIN produces the
better forecast in 5{,}386 of 7{,}657 decided country-months (70\%) yet loses
on aggregate: it wins small on the quiet mass and pays large on transitions,
the fingerprint of a system whose errors concentrate exactly where VIEWS's
covariates earn their keep. The horizon decomposition agrees: the deficit is
largest at h1--3 (.0413 vs .0337), where fresh signal matters most, and
smallest at h7--12 (.0503 vs .0458).
\end{enumerate}

We refrain from formal tests of Brier differences: the 7{,}680 records are
strongly dependent (overlapping 12-month horizons from five vantages;
within-country serial correlation), and a clustered Diebold--Mariano
treatment on five effective vantages would be underpowered theater. The
differences we interpret---pools versus singles, and the covariate nulls
below---are either zero-parameter comparisons or preregistered decisions,
not significance claims.

\section{The covariate program: five preregistered failures}\label{sec:covariates}

\subsection{The tune/validate protocol}\label{sec:protocol}

Walk-forward evaluation removes look-ahead but not \emph{designer
overfitting}: a researcher who iterates conditioning schemes against the
full scored history is fitting to it. After one such episode---a V-Dem
regime conditioning that looked plausible and measured worse on every
country suite (e.g., sb${\geq}25$ $.0411{\to}.0413$; all-type${\geq}100$
$.0366{\to}.0371$)---we froze a protocol:

\begin{enumerate}
\item Split scored years at 2007: \textbf{tune} (${\leq}2007$),
\textbf{validate} (${>}2007$).
\item Candidate schemes are enumerated \emph{a priori}; any thresholds are
computed from tune-era data only.
\item Candidates compete on tune; the winner is compared to baseline on
validate \textbf{exactly once}.
\item Adoption requires ${\geq}1\%$ relative Brier improvement on validate
\emph{and} a consistency guard on tune (a majority of cuts helping, for
parametric families; the winner itself beating baseline, for heterogeneous
sets).
\item Every consultation of validate is logged; when an era's validate set
has been consulted for a family of hypotheses, it is \textbf{retired}---
subsequent studies require genuinely fresh outcome years.
\end{enumerate}

The 1\% bar was set before any study ran, calibrated to the observed noise
scale of the validate split (a scheme that lost on tune still swung $+0.4\%$
on validate; see below). The protocol's output is a verdict, not a p-value;
its epistemic content is that validate data cannot be shopped.

\subsection{Results}

\paragraph{Study 1 --- youth bulges (country grain).} Descriptively, age
structure is the strongest single onset covariate we measured: among 6{,}280
currently-quiet country-years with covariate coverage (base onset rate
3.3\%), those with ${>}39\%$ of population under 14 onset at 6.5\% versus
1.7\% for older populations---a ${\sim}4\times$ lift, consistent with
\citet{urdal2006clash}. Under the protocol (candidate cuts at the
50th/60th/66th/75th tune-era percentiles): \textbf{zero of four cuts beat
the baseline engine on tune} (best .05384 vs .0537); the least-bad cut
carried a $+0.4\%$ relative edge to validate---below the bar, and with the
tune guard failed. \textbf{Rejected.} A naive adoption threshold
($\Delta\brier < -.0001$) would have adopted a covariate that could not beat
its own tuning baseline; the guard exists precisely for this case.

\paragraph{Study 2 --- dispute and alliance history (pair grain).} Four
candidates on the pair universe: any militarized-dispute history (ever-MID),
MID within 25 years, fatal-MID history, and a defense-pact flag (COW
alliances, last observed status carried forward). Descriptively the leading
candidate is enormous: among pair-years entering the \emph{cold} bucket,
those with MID history onset at 0.194\% (44/22{,}649) versus 0.006\%
(14/215{,}855) without---a \textbf{30$\times$} class-rate ratio. Under the
protocol: \textbf{zero of four candidates beat baseline on tune} (.000574 vs
.000573); the tune-selected candidate (ever-MID) was 0.7\% relatively
\emph{worse} on validate. \textbf{Rejected}---for a reason
\S\ref{sec:blindness} shows is structural rather than empirical.

\paragraph{Study 3 --- exclusion and the joint hypothesis (country grain;
final consultation of the 2007 era).} The one descriptively motivated
interaction not yet tested: ethnic exclusion alone, youth$\wedge$exclusion
(the compound cell onsets at 8.6\% versus 2.9\% for older/low-exclusion
years), and youth$\vee$exclusion---deliberately two-cell flags rather than
fragmenting cross-products. \textbf{Zero of three beat baseline on tune};
the tune-selected candidate (exclusion alone) was 1.0\% relatively worse on
validate---the worst held-out reading of any study. \textbf{Rejected}, and
the 2007 validation era was retired as preregistered (country validate
consulted twice, pair once).

Together with the two pre-protocol ablations (regime conditioning; horizon
decay), the scoreboard is \textbf{five families proposed, five rejected},
while the engine's own recency refinements (episode age, intensity, tempo,
neighbor) were adopted on the same walk-forward evidence. The pattern is not
that structural covariates are uninformative---descriptively they are
strong---but that they are informative \emph{about the same thing recency
already measures}. A young, exclusionary state that has been at peace for
thirty years is, empirically, at low risk; a middle-aged democracy one year
out of civil war is not. Conflict history is not a proxy to be improved
upon; it is the sufficient statistic the covariates approximate.

Two mechanisms recur. \textbf{Absorption}: conditional on bucket, the
covariate's class-rate difference shrinks toward zero (youth's $4\times$
marginal lift is mostly the young-countries-fight-recently composition
effect). \textbf{Fragmentation under shrinkage}: splitting a bucket cuts
class support; moment-matched shrinkage then pulls both children toward the
parent, so even a real difference buys little resolution while adding
estimator variance---the measured regime failure, and the reason the joint
study used two-cell flags. The third mechanism deserves its own subsection.

\subsection{Metric blindness at rare-event base rates}\label{sec:blindness}

The pair-grain rejection is \emph{not} evidence that dispute history is
uninformative---the $30\times$ class ratio is real and walk-forward. It is
evidence about the scoring rule. Decompose the maximum possible Brier
improvement from any refinement of a class with pooled rate $\bar p$ into
subclasses $c$ with masses $w_c$ and true rates $\pi_c$ (this is the
resolution term of Murphy's decomposition; \citealp{murphy1973new}):
\begin{equation}
\Delta\brier_{\max} \;=\; \sum_c w_c\,(\pi_c - \bar p)^2 .
\label{eq:blindness}
\end{equation}
For the cold-pair class (238{,}504 pair-years entering \emph{cold}):
$w_1 = 0.095$, $\pi_1 = .00194$; $w_0 = 0.905$, $\pi_0 = .000065$;
$\bar p = .000243$. Then $\Delta\brier_{\max} = 0.095 \cdot (.00170)^2 +
0.905 \cdot (.00018)^2 \approx 3.0 \times 10^{-7}$---\textbf{0.05\% of the
.000573 tune-era baseline}, twenty times below the preregistered 1\% bar. At
these base rates, \emph{perfect} conditioning is invisible to aggregate
Brier: the quiet mass is already scored nearly perfectly, and the flagged
class is too small for its (large, real) improvement to register. The same
arithmetic explains why the engine's $7\times$ neighbor-contagion signal
barely moves month-grain Brier, and why pair-grain skill ($+19.6\%$) looks
modest despite genuine discrimination.

The methodological point generalizes: \textbf{strictly proper scores
guarantee honest elicitation, not sensitivity to policy-relevant refinements
under extreme class imbalance.} A field that selects covariates by aggregate
Brier on rare-event targets will systematically discard its strongest onset
signals. Candidate remedies---weighted proper scores
\citep{ferro2011extremal}, the log score (which diverges precisely where
Brier saturates), or class-conditional evaluation---must be
\emph{preregistered}, because switching metrics after observing results
reintroduces the shopping the protocol exists to prevent. In our own system
the verdict stands (the era is spent); the $30\times$ signal ships as
displayed context on pair forecasts, labeled as having failed the adoption
bar. A log-loss bar for the pair grain is preregistered here for the next
evaluation era (outcome years ${\geq}2026$).

\section{Pooling: the best forecast is nobody's model}\label{sec:pools}

The head-to-head structure of \S\ref{sec:arena}---TOCSIN better on 70\% of
months, worse on aggregate---is the textbook precondition for combination
gains: the systems' errors are differently shaped. We therefore score
equal-weight linear opinion pools. Equal weights involve zero fitted
parameters, so scoring them on the same vantages is measurement, not tuning;
the adoption criterion (beat the best member) was fixed in advance, and all
four candidate pools are reported (Table~\ref{tab:pools}).

\begin{table}[t]
\centering
\caption{Equal-weight pools versus the best single model.}
\label{tab:pools}
\begin{tabular}{lrr}
\toprule
pool & Brier & vs best single \\
\midrule
VIEWS + persistence          & \textbf{.03987} & $\mathbf{-3.3\%}$ \\
VIEWS + persistence + TOCSIN & .04087 & $-0.9\%$ \\
VIEWS + TOCSIN               & .04144 & $+0.5\%$ \\
persistence + TOCSIN         & .04329 & $+5.0\%$ \\
\bottomrule
\end{tabular}
\end{table}

Two results. First, the field's cheapest available accuracy gain on this
target is not a model at all: averaging the ML system with naive persistence
beats everything, at every horizon band, for free. This is the
Bates--Granger result landing in the conflict domain with unusual force, and
it should discipline how single-model improvements are marketed: a new
system that does not beat the \emph{pool} has not advanced the frontier.

Second, the pool is a diagnostic instrument. Adding TOCSIN to
VIEWS$+$persistence \emph{worsens} it (.03987${\to}$.04087) despite TOCSIN's
70\% month-wise win rate, and TOCSIN$+$persistence is barely better than
persistence alone. The inference is sharp: at month grain, TOCSIN's forecast
is (approximately) a recalibrated persistence---its information is already
inside the pool, so it adds redundancy, not diversity. Pools pay for
\emph{orthogonal} error, and VIEWS's covariate machinery is the orthogonal
ingredient. This reframes the transparency tax of \S\ref{sec:arena}: the
transparent system's residual value on this target is not incremental
accuracy (persistence supplies its signal more cheaply) but
auditability---and its \emph{distinct} value lies on targets no other system
prices, to which we return below.

\section{Discussion}\label{sec:discussion}

\paragraph{A deflationary synthesis.} Assembled, the results argue that the
standard country-month occurrence target is close to its practical
predictability ceiling: (i) a pure recency machine reaches within 12\% of
the state of the art; (ii) naive persistence reaches within 5\%; (iii) five
theoretically central covariate families add nothing out-of-sample once
recency is conditioned on; (iv) the cheapest remaining gain is a
zero-parameter average. None of this says conflict is
unpredictable---skill versus climatology is $+50$--$57\%$ throughout---it
says the \emph{predictable part is mostly history}, and that the marginal
covariate, however strong its bivariate association, is purchasing
information the forecast already owns. This is the caution of
\citet{cederman2017predicting}, converted from warning into measurement.

\paragraph{What transparency bought.} The auditable design is not (on this
target) an accuracy strategy; it is an epistemics strategy. Because every
probability decomposes into named classes and counts, failures localize: we
can state \emph{why} youth fails (absorption), why regime fails
(fragmentation), why dispute history fails (metric blindness)---mechanisms,
not shrugs. The same architecture made the preregistration protocol cheap
enough to actually follow, and the negative results reportable at
publication grade. We suggest the field's forecasting systems adopt
adoption-protocols of this shape: enumerated candidates, tune/validate
separation, single consultation, retirement of spent eras, and published
rejections.

\paragraph{Rare events need preregistered metrics.} The decomposition of
\S\ref{sec:blindness} implies that aggregate Brier is the wrong adoption
criterion wherever base rates are far from the forecast probabilities'
working range---dyadic onset, interstate war, mass-atrocity onset. Weighted
or logarithmic criteria see what Brier cannot, but only preregistration
keeps them honest.

\paragraph{Limitations.} (1) The arena scores one target at one
grain---VIEWS's home turf; TOCSIN's design center (annual horizons,
arbitrary thresholds, termination, coups---the operational questions in its
journal) has no external comparator yet, and our claims are correspondingly
scoped. (2) Five vantages with overlapping horizons preclude meaningful
significance testing on Brier differences; our decisive comparisons are
zero-parameter or preregistered instead. (3) Seven percent of arena outcomes
were provisional (candidate GED) at scoring time; UCDP revisions could shift
third-decimal values. (4) VIEWS is represented by its published dichotomous
output under empirically verified semantics; other VIEWS surfaces (expected
fatalities) are not evaluated. (5) The covariate verdicts are specific to
this engine's conditioning structure and adoption bar; they bound the
\emph{increment over recency}, not the covariates' scientific relevance.
(6) The 2007-split era is spent; our own future covariate work must wait for
fresh outcome years, and we commit here to the preregistered log-loss bar
for the pair grain.

\paragraph{Future work.} The deployed system accumulates the missing
evidence automatically: a journal of operational forecasts resolves monthly
against UCDP updates, building both an annual-grain arena on TOCSIN's own
targets and the fresh validation era (outcomes ${\geq}2026$) that the next
covariate---a continuous tempo input aimed at the measured h1--3
deficit---must clear. The combination result suggests a standing published
pool alongside; the metric result suggests the community converge on
preregistered rare-event scoring before, not after, the next generation of
systems reports its gains.

\section{Reproducibility}\label{sec:repro}

Everything is public. Code (MIT) and the merged dataset (CC~BY~4.0, with
per-source attribution chains and codebook) are at
\url{https://github.com/KaraZajac/TOCSIN}; the live system, including the
arena table, all walk-forward calibration diagnostics, and every forecast
with its full reference-class ladder, is at
\url{https://tocsin.karazajac.io}. A single command (\texttt{tocsin pull \&\&
tocsin build \&\& tocsin backtest \&\& tocsin benchmark}) regenerates every
number in this paper from primary sources; the preregistration protocol is
\texttt{tocsin protocol}. The five-vantage arena pins the exact VIEWS runs
consumed. Raw ACLED data informs a display-only divergence indicator on the
site and is neither redistributed nor used in any result reported here.

\section*{Acknowledgments}

The system was designed, built, and audited by the author in an intensive
pair-collaboration with an AI assistant (Claude, Anthropic), which drafted
code and this manuscript under the author's direction; all design decisions,
verdicts, and errors are the author's. The author thanks the maintainers of
UCDP, VIEWS, the Correlates of War project, V-Dem, EPR, PRIO, Powell \&
Thyne, Gleditsch \& Ward, Our World in Data, and the World Bank for decades
of public data stewardship.

\bibliographystyle{apalike}
\bibliography{references}

\end{document}
